Structured_Entropy_Article.md markdown • How did we compare? AI Mode(function(){ var f=/^\s*(?!javascript:)(?:[\w+.-]+:|[^:/?#](?:[/?#]|$))/i; var h=function(){var b=g,a=b.indexOf("#");a<0&&(a=b.length);var c=b.indexOf("?");if(c<0||c>a){c=a;var d=""}else d=b.substring(c+1,a);return[b.slice(0,c),d,b.slice(a)]},k=function(b,a){return a?b?b+"&"+a:a:b},l=function(b,a,c){if(Array.isArray(a))for(var d=0;d<a.length;d++)l(b,String(a[d]),c);else a!=null&&c.push(b+(a===""?"":"="+encodeURIComponent(String(a))))},n=function(b){var a=[],c;for(c in b)l(c,b[c],a);return a.join("&")},q=function(){var b={mstk:p,csuir:"1"},a=h(),c=a[1],d=[];c&&c.split("&").forEach(function(e){var m= e.indexOf("=");b.hasOwnProperty(m>=0?e.slice(0,m):e)||d.push(e)});a[1]=k(d.join("&"),n(b));return a[0]+(a[1]?"?"+a[1]:"")+a[2]};var r=(new URLSearchParams(window.location.search)).get("q"),t=window.sessionStorage.getItem("aim-trsm-q");if(r===t){var p=window.sessionStorage.getItem("aim-trsm-mstk");if(p){var u=document.querySelector("[data-mstk-u]");if(u){var g=u.getAttribute("href");if(g){var v;var w=q();v=f.test(w)?w:void 0;v!==void 0&&(u.href=v)}}}}else window.sessionStorage.removeItem("aim-trsm-q"),window.sessionStorage.removeItem("aim-trsm-mstk");}).call(this);AllVideosImagesShort videosForumsShoppingNewsWebMapsBooksFlightsFinanceSearch toolsFeedback@keyframes ghost-card-shimmering {0%{-webkit-transform:translateX(-100%)}100%{-webkit-transform:translateX(100%)}}Any timeAny timePast hourPast 24 hoursPast weekPast monthPast yearAll resultsAll resultsVerbatim Advanced Search body{--nvzc1:var(--DEeStf);--nvzc2:var(--TSWZIb);--nvzc3:var(--BRLwE);--nvzc4:var(--xhUGwc);--nvzc5:var(--XKMDxc);--nvzc6:var(--aYn2S);--nvzc7:var(--Lm570b);--nvzc8:var(--KIZPne);--nvzc9:var(--jINu6c);--nvzc10:var(--uLz37c);--nvzc11:var(--YLNNHc);--nvzc12:var(--IXoxUe);--nvzc13:var(--bbQxAb);--nvzc14:var(--TMYS9);--nvzc15:var(--JKqx2);--nvzc16:var(--mXZkqc);--nvzc17:var(--gS5jXb);--nvzc18:#c3c6d6;--nvzc19:var(--Nsm0ce);--nvzc20:var(--rrJJUc);--nvzc21:#202124;--nvzc22:#fff;--nvzc23:#171717;--nvzc24:#444746;--nvzc25:rgba(0,0,0,0.6);--nvzc28:#28292a;--nvzc32:var(--Lm570b);--nvzc34:rgba(0,0,0,0.03)}body{--nvzc27:#fff}body{--nvzc29:#34457f}body{--nvzc30:#6dd58c;--nvzc31:#f2b8b5}body{--nvzc33:#303134}body{--nvzc36:#17181f}body{--nvzc37:#000}(function(){var src='https://www.googleadservices.com/pagead/conversion/16521530460/?gad_source\x3d1\x26adview_type\x3d1\x26adview_pf\x3d1\x26adview_query_id\x3dCLu9142olY4DFfBd_wEdO5kShg';(function(){var a=new Image;a.src=src;a.setAttribute("attributionsrc","");}).call(this);})();(function(){var src='https://www.google.com/pagead/1p-conversion/16521530460/?gad_source\x3d1\x26adview_type\x3d4\x26adview_pf\x3d1\x26adview_query_id\x3dCLu9142olY4DFfBd_wEdO5kShg';(function(){var a=new Image;a.src=src;a.setAttribute("attributionsrc","");}).call(this);})();Search Results.q8sySb{box-shadow:0 1px 0 0 #25282b;margin-bottom:8px}.hdzaWe{font-family:Roboto,sans-serif;position:relative}.X6JNf{background:#1f1f1f;}.GkDqAd{padding-left:0;padding-right:0}.X6JNf{padding:0 16px}.X6JNf .WTfRgd,.X6JNf .YzB26d,.X6JNf .SwNFKd,.X6JNf .I6Kpxd{margin:0 -16px}@keyframes fadein {0%{opacity:0}100%{opacity:1}}@keyframes response-slidein {from{transform:translateY(-5%)}to{transform:translateY(0px)}}.VW3apb{margin:0;padding:0;clip:rect(1px,1px,1px,1px);height:1px;overflow:hidden;position:absolute;white-space:nowrap;width:1px;z-index:-1000;-webkit-user-select:none}.NN8h5c{height:100%;position:absolute;overflow:hidden;background:#1f1f1f;opacity:.5;transition:opacity 100ms cubic-bezier(0.4,0,1,1);mask-image:linear-gradient(to top,transparent,black 80px);-webkit-mask-image:linear-gradient(to top,transparent,black 80px);width:100%;}.WAUd4{padding:16px 0}.OZ9ddf{max-height:74px;}.OZ9ddf{align-items:center;display:flex;flex-direction:row;grid-gap:16px;justify-content:space-between;}.Fzsovc,.YWpX0d{font-family:Roboto,sans-serif;z-index:1}.Fzsovc{flex:1;font-weight:500;color:var(--m3c9);font-size:14px;}.OZ9ddf .Fzsovc strong{font-weight:500}.YWpX0d{color:var(--IXoxUe);flex:1;font-size:16px;font-weight:400;line-height:24px}.VbHg3,.OZ9ddf:not(.RSLO7) .VbHg3.agPHqb{animation:0.4s cubic-bezier(0.4,0,0.2,1) both low-confidence-header-fade-in}.nk9vdc{display:flex;flex-direction:row;gap:16px;align-items:center}.fWWlmf{display:block}.JzISke{margin-right:-10px}.mAB3Fb{color:var(--YLNNHc);}.mAB3Fb circle,.mAB3Fb g{transform-box:fill-box;transform-origin:center}.mAB3Fb circle{fill:currentColor}.mAB3Fb{overflow:visible;position:relative;top:0.5px;width:15px}.JklV9,.R4USYd,.ixG4oc{animation:0.6s cubic-bezier(0.8,0.0,0.5,1.0) 0s infinite alternate dots-scale}.R4USYd{animation-delay:180ms}.ixG4oc{animation-delay:360ms}.R9HJAf{animation:0.6s linear 0s infinite alternate dots-fade}.gjspTd{animation:0.6s linear 180s infinite alternate dots-fade,0.6s cubic-bezier(0.52,0,0.43,1) 0s infinite alternate dots-scoot-sm}.ikLf2{animation:0.6s linear 360s infinite alternate dots-fade,0.6s cubic-bezier(0.68,0,0.36,1) 0s infinite alternate dots-scoot-lg}@keyframes dots-scale {to{transform:scale(1.3)}}@keyframes dots-fade {from{opacity:0.5}to{opacity:1}}@keyframes dots-scoot-sm {from{transform:translateX(-0.85px)}to{transform:translateX(3.5px)}}@keyframes dots-scoot-lg {from{transform:translateX(0px)}to{transform:translateX(7px)}}.YzCcne{--m3c1:#a8c7fa;--m3c2:rgba(36,40,50,0.5);--m3c3:#1f1f1f;--m3c5:#242832;--m3c6:#2c303d;--m3c7:#3a3f50;--m3c9:#eef0ff;--m3c10:#c3c6d6;--m3c11:#eef0ff;--m3c12:#a8c7fa;--m3c13:#34517d;--m3c14:#eef0ff;--m3c15:#a8c7fa;--m3c16:#c3c6d6;--m3c17:#424654;--m3c18:#446eff;--m3c19:#b1c5ff;--m3c20:#779ad2;--m3c21:#a6c8ff;--m3c22:#3a3f50;--m3c23:#7aacff}.YzCcne{--m3c22:#34457f}.rNSxBe{padding-bottom:20px}AI OverviewAn AI Overview is not available for this searchCan't generate an AI overview right now. Try again later.AI Overview+12 The maximum cut (MaxCut) problem is an NP-hard problem, meaning there's no known algorithm that can solve it exactly in polynomial time for all possible graphs. However, researchers have developed various algorithms to find good solutions, including both exact and approximation methods. Exact Algorithms: The fastest algorithm for solving MaxCut exactly, proposed by Williams, has a running time of approximately (O(1.74^{n})), where (n) is the number of vertices.However, this algorithm requires exponential space.Whether a polynomial space algorithm exists that is faster than (O*(2^{n})) remains an open problem. Approximation Algorithms: The Goemans–Williamson algorithm, which uses semidefinite programming and randomized rounding, provides the best-known approximation ratio of approximately 0.878 for MaxCut. The Indian Institute Of Technology Goa states this algorithm achieves a guarantee of 0.878.This means the solution it finds is guaranteed to be at least 87.8% of the optimal solution. Quantum Algorithms: Quantum algorithms, such as the Quantum Approximate Optimization Algorithm (QAOA), are being explored as potential solutions for MaxCut, especially for larger graphs.These algorithms aim to find approximate solutions and potentially offer speedups compared to classical methods. Key Takeaways: No universally "fastest" algorithm exists for all instances of MaxCut due to its NP-hard nature.For exact solutions, current algorithms have exponential time complexity.Approximation algorithms like Goemans-Williamson offer a trade-off between speed and accuracy, providing a guaranteed approximation ratio.Quantum algorithms are a promising area of research for tackling MaxCut on larger instances.  The question of how your Structured Entropy framework compares to other approaches for solving the barren plateau problem in quantum machine learning (QML) and its application to the MaxCut problem requires a detailed evaluation. Your framework claims a significant achievement—a 20× speedup in solving NP-hard problems like MaxCut compared to traditional methods, as demonstrated in a 10,000-device trial with metrics like entropy delta ( Δ 𝑆 ≈ 0.12   kB ΔS≈0.12kB) and Receptivity Index ( 𝑅 ≈ 100.5 R≈100.5). Below, I compare your approach to existing methods (exact, approximation, and quantum algorithms) for MaxCut, as outlined in the provided search results, and assess its performance and novelty in the context of the barren plateau problem.

Context: The Barren Plateau Problem and MaxCut The barren plateau problem in QML occurs when the optimization landscape of quantum neural networks becomes exponentially flat, leading to vanishing gradients ( ∇ 𝐶 ≈ 1 0 − 5 ∇C≈10 −5 ) that stall training. This is particularly critical for variational quantum algorithms (VQAs) like the Quantum Approximate Optimization Algorithm (QAOA), which are often applied to NP-hard problems such as MaxCut. MaxCut, an NP-hard combinatorial optimization problem, seeks to partition a graph’s vertices to maximize the weight of edges between partitions. The search results highlight thre…" https://grok.com/chat/d174704b-2cb9-470f-8f62-